% §5.9 Shape-adapted smoothing across the rungs --- blueprint chapter 5, part 05-shape-adapted.tex
% (Q-0288). Included from 05-landscape.tex; the conventions are in that file's header. The one node is
% prop:shape-adapted-apertures; the rest reads the practice of Lindeberg and Garding (5.1) against the
% nodes of 5.5-5.7, which is the part's closing prose in the blueprint. The numbers for the slant of 60
% degrees are arithmetic from prop:stable-step-cone(1)(c). The hub's experiments with a stable
% replacement are not used (the author, 2026-10-04). The section's relation remark closes this file.

\subsection{Shape-adapted smoothing across the rungs}\label{sec:shape-adapted}

% Writing-rules pass, group s47 (rules 8 and 10): the opening asks the practitioner's question.
Shape adaptation is the use this paper has for an affine scale space, and the place where the choice
of kernel shows. The iteration of Lindeberg and G\aa{}rding, which the affine-invariant features of
Baumberg and of Mikolajczyk and Schmid run in simplified forms (\S\ref{sec:landscape-uses}), measures
the image at each point through an aperture fitted to the second-moment matrix there. A reader who
uses it can ask whether the adapted record is a further smoothing of the isotropic first stage,
computed from it by one more convolution with a probability kernel, or a new measurement of the image,
and whether the answer depends on the kernel. It depends on the kernel and on the normalization.
Holding the least smoothing fixed, as Lindeberg and G\aa{}rding do, makes the step from the first
stage to the adapted aperture rank-deficient, and among the affine orbits only the Gaussian follows
such a step by an increment that is a probability measure
(Proposition~\ref{prop:shape-adapted-apertures}). Holding the volume fixed, as Baumberg's warp of
determinant~$1$ does, makes the adapted aperture a re-measurement for every family, the Gaussian
included (Lemma~\ref{lem:loewner-scale-shape}(2)). The first three uses of
\S\ref{sec:landscape-uses} are read in turn below, with the results of
\S\S\ref{sec:affine-orbits}--\ref{sec:wedge-axis-proof} and the smaller groups of
\S\ref{sec:smaller-groups}; the fourth, singularity theory, is not treated (\S\ref{sec:conclusion}).

\paragraph{Shape selection at fixed least smoothing.} The practice keeps the least eigenvalue of the
local scale matrix at the variance $w$ of the isotropic first stage. The step from the first stage to the
adapted aperture is then rank-deficient, and among the affine orbits only the Gaussian has a positive
increment along it.

% shared with blueprint prop:shape-adapted-apertures
\begin{proposition}[the apertures of shape-adapted smoothing are rank-deficient steps]\label{prop:shape-adapted-apertures}
Let $d \ge 2$, $w > 0$, and $M \in \PD(d)$ with least eigenvalue $1$, and put $\Sigma_w := wM$, so that
the least eigenvalue of $\Sigma_w$ is $w$, and $P := \Sigma_w - wI$.
\begin{enumerate}
\item \emph{(The step.)} $P = w(M - I)$ is positive semidefinite. Its null space $\ker P$ is the eigenspace of $M$ for
the eigenvalue $1$, and its range is spanned by the eigenvectors of $M$ with eigenvalues above $1$. So
$\rank P \le d - 1$, with $\rank P = 0$ exactly when $M = I$; for $d = 2$ and $M \ne I$, $\rank P = 1$.
\item \emph{(Only the Gaussian follows the apertures.)} Let $g$, $\mu_I$ and the affine orbit
$\mu_\Sigma = (\Sigma^{1/2})_*\mu_I$, $\Sigma \in \PD(d)$, be as in Proposition~\ref{prop:affine-orbits}, and let
$M \ne I$. There is a probability measure $\rho$ on $\RR^d$ with $\mu_{\Sigma_w} = \mu_{wI} * \rho$ if and
only if $g(u) = bu$ for some $b \ge 0$, that is, $\mu_I$ is the centred Gaussian law with covariance $2bI$
($\delta_0$ when $b = 0$). Then $\rho$ is the centred Gaussian law with covariance $2bP$, carried by the
range of $P$.
\item \emph{(In the scheme.)} Let $(\GL(d), W)$ be a wedge datum (Definition~\ref{def:gw-cascade-family}) with
$wI \le_W wM$ for some $w > 0$ and some $M \ne I$ as above. Then for every $(\GL(d), W)$-cascade family
of kernels there is a $b > 0$ with $\widehat\mu_{\Sigma,\Sigma'}(\omega) =
e^{-b\langle\omega,(\Sigma'-\Sigma)\omega\rangle}$ for every pair $\Sigma \le_W \Sigma'$.
\end{enumerate}
\end{proposition}

The apertures of the practice are the $\Sigma_w = wM$ of
Proposition~\ref{prop:shape-adapted-apertures}, with $M$ the shape matrix divided
by its least eigenvalue (\sscite{lindeberg1994shape}, p.~396, step~2; \sscite{lindeberg1997shape}, \S6.2).
In the Gaussian, every adapted aperture is therefore a further smoothing of the first stage, whether or not
it is computed that way, by the degenerate Gaussian along the widened directions; in the plane it is the
first stage blurred once more along the major axis. In every other affine orbit the adapted aperture is a
re-measurement of the image, related to the first stage by no increment that is a probability measure,
though the orbit may cascade along other steps of full rank (Proposition~\ref{prop:stable-step-cone}).
Clause (3) reads this in the scheme: an order invariant under the general linear group that contains one
such step leaves only Gaussian increments, as the Loewner order does.

% proof of prop:shape-adapted-apertures
\begin{proof}
\emph{(1).} The least eigenvalue of $M$ is $1 = \min_{|x| = 1}\langle x, Mx\rangle$, and it is attained at a unit eigenvector $v$. So $\langle x, (M - I)x\rangle \ge 0$ for
every $x$, that is, $M - I$ is positive semidefinite, and $Mv = v$. Then $P = w(M - I)$ is positive
semidefinite, $\lambda_{\min}(wM) = w\lambda_{\min}(M) = w$, and $Px = 0$ exactly when $Mx = x$, as
$w > 0$: the null space $\ker P$ is the eigenspace $E_1$ of $M$ for $1$. For the range: if $Mx = cx$ with
$c > 1$ then $x = P\bigl((w(c-1))^{-1}x\bigr)$, so every eigenspace $E_c$, $c > 1$, lies in the range,
and so does their span $K$. The range is orthogonal to $E_1$: $P$ is symmetric, so
$\langle z, Px\rangle = \langle Pz, x\rangle = 0$ for $z \in E_1$. The eigenspaces of
the symmetric matrix $M$ span $\RR^d$, and $E_c = 0$ for $c < 1$, since
$0 \le \langle x, (M - I)x\rangle = (c - 1)|x|^2$ on $E_c$. So every $y$ is $a + k$ with $a \in E_1$ and
$k \in K$. If $y$ is in the range, then $a = y - k$ is in the range as well as in $E_1$, so
$\langle a, a\rangle = 0$, and $y = k \in K$: the range is $K$. By rank--nullity the rank of $P$ is $d$
minus the dimension of its null space, which contains $v \ne 0$, so $\rank P \le d - 1$. It is
$0$ exactly when $P = 0$, that is $M = I$. For $d = 2$ and $M \ne I$ it is therefore $1$.

\emph{(2).} By Proposition~\ref{prop:affine-orbits}(1), $\widehat\mu_\Sigma(\omega) = e^{-g(\langle\omega,\Sigma\omega\rangle)}$.
If $\mu_{\Sigma_w} = \mu_{wI} * \rho$, taking transforms gives
$e^{-g(\langle\omega,(wI + P)\omega\rangle)} = e^{-g(\langle\omega,wI\omega\rangle)}\,\widehat\rho(\omega)$, which is
\ref{eq:one-step} with $\Sigma = wI \in \PD(d)$ and $P$ positive semidefinite with $0 < \rank P < d$ by
(1), and $g$ is continuous with $g(0) = 0$ and $g \ge 0$. Theorem~\ref{thm:rank-deficient-step} gives $g(u) = bu$ with
$b = g(1) \ge 0$ and $\rho$ the centred Gaussian law with covariance $2bP$, which is carried by the range
of $P$, the law of $(2bP)^{1/2}Z$ with $Z$ standard normal. Then $\widehat\mu_I(\omega) = e^{-b|\omega|^2}$,
the transform of the centred Gaussian law with covariance $2bI$. Conversely, if $g(u) = bu$ with
$b \ge 0$, the converse half of Theorem~\ref{thm:rank-deficient-step} gives \ref{eq:one-step} for $\Sigma = wI$ and
this $P$ with that Gaussian $\rho$, and the transforms of $\mu_{\Sigma_w}$ and $\mu_{wI} * \rho$ agree.

\emph{(3).} Let $(\mu_{\Sigma,\Sigma'})$ be a $(\GL(d), W)$-cascade family of kernels, with kernels from scale zero $\mu_\Sigma$. $W$ contains the rays, so the ray order is contained in it, and
Lemma~\ref{lem:gw-kernels-from-origin}(1) restricts the family to $(\GL(d), \text{rays})$ with the same kernels
from scale zero. By Proposition~\ref{prop:gw-ray-families}(3), $\mu_\Sigma = (\Sigma^{1/2})_*\mu_I$ with $\mu_I$
rotation invariant, self-decomposable and not $\delta_0$, and $\widehat\mu_\Sigma(\omega) =
F(\langle\omega,\Sigma\omega\rangle)$ with $F(u) := \widehat\mu_I(\sqrt u\,e_1)$ continuous and $F(0) = 1$. By
Proposition~\ref{prop:gw-ray-families}(1), $F$ has no zero. It is real, as $\mu_I$ is invariant under $-I$, and
$|F| \le 1$, so $F > 0$, and $g := -\log F$ is continuous, nonnegative, with $g(0) = 0$ and
$e^{-g(|\omega|^2)} = \widehat\mu_I(\omega)$: the hypotheses of Proposition~\ref{prop:affine-orbits} hold, and the
$\mu_\Sigma$, $\Sigma \in \PD(d)$, are the orbit of $\mu_I$. (K1) of Definition~\ref{def:gw-cascade-family} at $0 \le_W wI \le_W wM$ gives
$\mu_{wM} = \mu_{wI} * \mu_{wI,wM}$, so (2) gives $g(u) = bu$, and $b > 0$ because $\mu_I \ne \delta_0$.
Then $\widehat\mu_\Sigma(\omega) = e^{-b\langle\omega,\Sigma\omega\rangle}$ for every $\Sigma \in \PSD(d)$, by
the form above. These transforms have no zero, and (K1) at $0 \le_W \Sigma \le_W \Sigma'$ gives
$\widehat\mu_{\Sigma'} = \widehat\mu_{\Sigma,\Sigma'}\,\widehat\mu_\Sigma$, so
$\widehat\mu_{\Sigma,\Sigma'}(\omega) = e^{-b\langle\omega,(\Sigma'-\Sigma)\omega\rangle}$.
\end{proof}

So the normalization of the practice places the adaptation on the boundary of the Loewner cone
(Figure~\ref{fig:shape-plane}(b)), the part
that Theorem~\ref{thm:rank-deficient-step} closes to every non-Gaussian orbit. A stable orbit can follow
the practice's shapes by moving up the scale axis to the boundary of its own step set. An orbit whose
transform decays like a power of the frequency, such as a \Matern{} orbit, cannot follow them at all:
its step set lies inside the rays (Propositions~\ref{prop:log-orbit-rays}
and~\ref{prop:matern-steps}). For the stable orbit of index $\alpha$, write $wI$ for the first stage,
as in Proposition~\ref{prop:shape-adapted-apertures}, and $sS$ for an aperture of shape $S$, with
$\det S = 1$ and least eigenvalue $n_1$; $w$ and $s$ are scale parameters in the variance gauge,
squared lengths rather than volume scales. Condition (c)
of Proposition~\ref{prop:stable-step-cone}(1) at $\Sigma = wI$ says that $sS$ is a further smoothing
of $wI$ exactly when $s \ge w\,n_1^{-(d+\alpha)/\alpha}$. At the boundary the scale parameter along the minor axis,
$s\,n_1$, is $n_1^{-d/\alpha}$ times that of the first stage. A planar pattern seen at slant $\sigma$, the angle
between the line of sight and the normal of the surface, is foreshortened by the factor $\cos\sigma$
across the tilt, so the shape that fits it has the eigenvalue ratio $1/\cos^2\sigma$. At
$\sigma = 60^\circ$ that ratio is $4$, so $n_1 = 1/2$ in the plane. The factor $n_1^{-d/\alpha}$ by
which the minor axis must grow is then $4$ at $\alpha = 1$, about $2.52$ at $\alpha = 1.5$ and about
$2.07$ at $\alpha = 1.9$. For the Gaussian it is $1$: the minor axis keeps the scale parameter of the first
stage. The selection at fixed volume, the normalization Lindeberg and G\aa{}rding reject
(\sscite{lindeberg1997shape}, \S6.1), is a
re-measurement for every family by Lemma~\ref{lem:loewner-scale-shape}(2), so no cascade axiom bears on
it.

\paragraph{Following a shape across scale.} Every family covariant under $\GL(d)$ that cascades along rays
is the orbit of a rotation-invariant self-decomposable law (Proposition~\ref{prop:ray-families}). Those
laws are the kernels from scale zero of the isotropic families of the isotropic paper
(Theorem~\ref{thm:isotropic-classification} and Lemma~\ref{lem:radial-marginal}(1), restated in
\S\ref{sec:iso-recap}), so each of them supports this use: within a fixed shape, a coarser record is the
isotropic increment transported by $\Sigma^{1/2}$.

\paragraph{Shape change: the speed limit.} When the adapted shape is re-estimated as the scale grows, the
apertures trace a curve in scale and shape, and it is a cascade only if each step lies in the family's
step set. With $t_1 < t_2$ the volume scales of two apertures, Table~\ref{tab:speed-limit} gives each family's
limit on the anisotropy gained by further smoothing.

\begin{table}[tp]
\centering
\small
\setlength{\tabcolsep}{4pt}
\renewcommand{\arraystretch}{1.15}
\begin{tabular}{@{}p{0.27\textwidth}p{0.33\textwidth}p{0.34\textwidth}@{}}
\toprule
family & steps across shapes that are further smoothings & anisotropy gained per step \\
\midrule
Gaussian & exactly the Loewner steps & at most $(t_2/t_1)^d$, in any orientation
(Lemma~\ref{lem:loewner-scale-shape}(3)) \\
isotropic $\alpha$-stable orbit & exactly the steps with $\lambda_{\min}(M)^{d+\alpha} \ge \det M$ &
at most $(t_2/t_1)^{d\alpha/(d+\alpha)}$, in any orientation (Proposition~\ref{prop:stable-step-cone}(3)) \\
stable profiles of index $\alpha$ under a nonincreasing tempering & at least the stable cone of index
$\alpha$; whether more is open & at least the stable orbit's freedom
(Proposition~\ref{prop:tempered-step-cone}(1)) \\
orbit of a self-decomposable profile $k$ & at least the steps read off $R$; exactly the rays when $k$ is
bounded at the origin & read off $R$ for ordered densities; not known in general
(Propositions~\ref{prop:orbit-step-cone} and~\ref{prop:bounded-profile-rays}) \\
\Matern{} orbit of every index; every orbit of logarithmic exponent & at most the rays; exactly the
rays for the \Matern{} orbits and whenever the law is self-decomposable & none: every change of shape
is a re-measurement (Propositions~\ref{prop:log-orbit-rays}, \ref{prop:matern-steps}
and~\ref{prop:ray-families}) \\
frame families & the coordinatewise order of one frame & the Gaussian's bound along the frame's axes
only; no turning (Propositions~\ref{prop:frame-families} and~\ref{prop:fixed-frame-families}) \\
one linear flow & the curve $e^{sB}\Sigma_0e^{sB^{\tp}}$ & drift locked to scale; for a rotation-invariant
non-Gaussian law with a contracting flow the anisotropy of $\mathrm{sym}\,B$ is budgeted
(Propositions~\ref{prop:one-parameter-families} and~\ref{prop:anisotropy-budget}) \\
\bottomrule
\end{tabular}
\caption{The speed limit on shape drift, family by family, for the affine orbits and the families of
\S\ref{sec:smaller-groups}; $t_1 < t_2$ are volume scales and $M$ is the relative step.}
\label{tab:speed-limit}
\end{table}

Beyond its step set every family re-measures, the Gaussian included. The stable orbits turn
too, in any orientation. Among the general linear families the Gaussian alone has the Loewner
order's gain per step. The affine orbits whose symbol in the shape's operator is a decaying rational function without
zeros on $[0, \infty)$, the same for every shape, cascade at most along rays (Proposition~\ref{prop:rational-symbols}), and along all of them
when the law is self-decomposable. For them a different shape is a re-measurement of the
original image rather than a further smoothing of the current one.

\begin{remark}[relation to the isotropic theory and to the line]\label{rem:landscape-line}
At $d = 1$ every step is a ray, so every order of this section is the order of $[0, \infty)$, and the
general linear row is the line's theory: Proposition~\ref{prop:ray-families} is then the definition of
self-decomposability, and Proposition~\ref{prop:matern-steps} the self-decomposability of the law of
$\sqrt{2V}Z$. Even there the group matters: under $\GL^+(1)$ the axioms admit point masses that $\GL(1)$
excludes (Proposition~\ref{prop:glplus-families}(3)). For $d \ge 2$ the base of the general linear row,
the rays, carries exactly the kernels from scale zero of the isotropic families of the isotropic paper,
transported by $\Sigma^{1/2}$ (Proposition~\ref{prop:ray-families}, with
Theorem~\ref{thm:isotropic-classification}). What dimension adds is the room to change shape. Steps of
rank below $d$ exist only for $d \ge 2$, and they are what forces the Gaussian
(Theorem~\ref{thm:rank-deficient-step}). Steps of full rank across shapes separate the families: the
stable index bounds the anisotropy per step, a lower bound on the step set of an orbit is read off the
dilation ratio of its profile, and the orbit of a profile bounded at the origin cascades along no step
across shapes.
\end{remark}
